Um operador é um símbolo que indica a realização de uma operação entre 
valores de um programa. Na linguagem C é possível encontrar alguns 
diferentes tipos de operadores, dentre eles:

\subsubsection{Aritméticos}

Os operadores aritméticos são utilizados para representar 
operações matemáticas entre valores. Segue abaixo uma lista de todos os 
operadores aritméticos:

\begin{center}
    \begin{tabular}{|c|c|c|c|}
        \hline
        Operador & Descrição & Exemplo & Resultado \\
        \hline
        + & soma & 10 + 20 & 30 \\
        \hline
        - & subtração & 20 - 10 & 10 \\
        \hline
        * & multiplicação & 10 * 2 & 20 \\
        \hline
        / & divisão & 20 / 10 & 2 \\
        \hline
        \% & módulo & 21 \% 10 & 1 \\
        \hline
        \texttt{++} & incremento & \texttt{x = 10; ++x} & 11 \\
        \hline
        \texttt{--} & decremento & \texttt{x = 10; --x} & 9 \\
        \hline
    \end{tabular}
\end{center}

Os operadores \texttt{++} e \texttt{--} exigem uma variável ou objeto
modificável. Em \texttt{++x}, o valor já incrementado é retornado, enquanto
\texttt{x++} retorna o valor anterior ao incremento. Quando usados
isoladamente, ambos deixam a variável com o mesmo valor final.

\begin{excode}
    int x = 10;
    int a = ++x; // x = 11, a = 11

    int y = 10;
    int b = y++; // y = 11, b = 10
\end{excode}

\begin{postscript}
    O operador de módulo realiza uma divisão entre os dois valores e retorna o 
    resto da operação.
\end{postscript}

\subsubsection{Relacionais}

Os operadores relacionais possuem a finalidade de fazer comparações entre 
valores, como por exemplo, verificar se um número é maior que outro. O 
resultado de uma comparação sempre vai ser \texttt{verdadeiro} ou 
\texttt{falso}.

Segue abaixo uma lista dos comparadores relacionais:

\begin{center}
    \begin{tabular}{|c|c|c|c|}
        \hline
        Operador & Descrição & Exemplo & Resultado \\
        \hline
        == & igual & 10 == 10 & verdadeiro \\
        \hline
        != & diferente & 10 != 10 & falso \\
        \hline
        \textgreater{} & maior que & 10 \textgreater{} 10 & falso \\
        \hline
        \textgreater{}= & maior ou igual & 10 \textgreater{}= 10 & verdadeiro \\
        \hline
        \textless{} & menor que & 10 \textless{} 10 & falso\\
        \hline
        \textless{}= & menor ou igual & 10 \textless{}= 10 & verdadeiro \\
        \hline
    \end{tabular}
\end{center}

\subsubsection{Lógicos}

Os operadores lógicos são usados em conjunto com os operadores relacionais. 
Para compreendê-los corretamente, é importante entender que na linguagem C 
os valores de \texttt{verdadeiro} e \texttt{falso} podem ser representados 
pelos números 1 e 0, respectivamente.

Ao utilizar os operadores lógicos, podemos criar verificações de valores 
mais complexas. Abaixo está a lista dos operadores existentes e como usá-los:

\begin{center}
    \begin{tabular}{|c|c|c|c|}
        \hline
        Operador & Descrição & Exemplo & Resultado \\
        \hline
        \&\& & AND lógico & 1 \&\& 1 & 1 \\
        \hline
        \textbar{}\textbar{} & OR lógico & 1 \textbar{}\textbar{} 0 & 1 \\
        \hline
        ! & NOT lógico & !1 & 0 \\
        \hline
    \end{tabular}
\end{center}

\begin{postscript}
    Caso tenha ficado um pouco confuso como cada operador atua exatamente, vou tentar 
    detalhar de maneira mais sucinta e com mais exemplos de como cada operador deve agir 
    em cada caso.

    \begin{enumerate}
        \item{\textbf{AND lógico:} Retorna verdadeiro apenas se ambos os 
            valores forem verdadeiros.}

            \begin{center}
                \begin{tabular}{|c|c|}
                    \hline
                    Exemplo & Resultado \\
                    \hline
                    0 \&\& 0 & 0 \\
                    \hline
                    0 \&\& 1 & 0 \\
                    \hline
                    1 \&\& 0 & 0 \\
                    \hline
                    1 \&\& 1 & 1 \\
                    \hline
                \end{tabular}
            \end{center}

        \item{\textbf{OR lógico:} Retorna verdadeiro se pelo menos um dos 
            valores for verdadeiro.}

            \begin{center}
                \begin{tabular}{|c|c|}
                    \hline
                    Exemplo & Resultado \\
                    \hline
                    0 \textbar{}\textbar{} 0 & 0 \\
                    \hline
                    0 \textbar{}\textbar{} 1 & 1 \\
                    \hline
                    1 \textbar{}\textbar{} 0 & 1 \\
                    \hline
                    1 \textbar{}\textbar{} 1 & 1 \\
                    \hline
                \end{tabular}
            \end{center}

        \item{\textbf{NOT lógico:} Inverte o valor lógico; verdadeiro se 
            torna falso e vice-versa.}

            \begin{center}
                \begin{tabular}{|c|c|}
                    \hline
                    Exemplo & Resultado \\
                    \hline
                    !0 & 1 \\
                    \hline
                    !1 & 0 \\
                    \hline
                \end{tabular}
            \end{center}
    \end{enumerate}
\end{postscript}

\subsubsection{Atribuição}

Os operadores de atribuição tem a função de atribuir valores a variáveis. Por exemplo, 
utilizando um operador de atribuição poderíamos dizer que A = 10 então toda vez que 
utilizássemos o valor A no código de um programa, ele seria interpretado como o número 10.

\begin{center}
    \begin{tabular}{|c|c|c|c|}
        \hline
        Operador & Descrição & Exemplo & Equivale \\
        \hline
        = & Atribuição de valor a variável & A = 10 & A = 10 \\
        \hline
        += & Atribuição com soma & A += 10 & A = A + 10 \\
        \hline
        -= & Atribuição com subtração & A -= 10 & A = A - 10 \\
        \hline
        *= & Atribuição com multiplicação & A *= 10 & A = A * 10 \\
        \hline
        /= & Atribuição com divisão & A /= 10 & A = A / 10 \\
        \hline
        \%= & Atribuição com resto & A \%= 10 & A = A \% 10 \\
        \hline
    \end{tabular}
\end{center}
